% !TeX root = ../navier-stokes-manuscript.tex
\subsection{Velocity, acceleration, jerk, and the Tower}\label{sub:BodyFiniteTimeResponseEscalation}

The finite endpoint has already forced the same velocity history through
unbounded critical levels.  Naming those levels does not move the fluid
between them.  On every strict subinterval that movement is carried by the
fundamental theorem of calculus.  Let \(X\) be one fixed spatial Banach space
in which the classical responses below are defined.  For every \(n\geq0\),
\begin{equation}\label{eq:BodyResponseIncrement}
  \norm{\partial_t^nu(t_1)-\partial_t^nu(t_0)}_X
  \leq
  \int_{t_0}^{t_1}\norm{\partial_t^{n+1}u(s)}_X\,ds.
\end{equation}
If \(\partial_t^{n+1}u\) had a uniform roof before a finite \(T_*\), its
integral would give \(\partial_t^nu\) a uniform roof as well.  Reading the
contrapositive upward, under the fixed-space and absolute-continuity hypotheses
spelled out in \cref{sub:FiniteTimeResponseEscalation},
\begin{equation}\label{eq:BodyConditionalTemporalEscape}
  \sup_{0<t<T_*}\norm{u(t)}_X=\infty
  \quad\Longrightarrow\quad
  \sup_{0<t<T_*}\norm{\partial_t^nu(t)}_X=\infty
  \quad\text{for every finite }n.
\end{equation}
The critical escape in \eqref{eq:BodyFiniteEndpointForcesCriticalEscape}
supplies the premise with \(X=\dot H^{1/2}(\R^3)\).  It does not prove that a
finite endpoint exists.  It says that if the history ends there, no finite
Eulerian time row may remain uniformly calm while the lower row escapes:
\[
  u,
  \qquad
  \partial_tu,
  \qquad
  \partial_t^2u,
  \qquad
  \partial_t^3u,
  \qquad
  \ldots
\]
Write
\begin{equation}\label{eq:TemporalTowerDefinitionBody}
  V_n:=\partial_t^nu,
  \qquad
  Q_n:=\partial_t^np,
  \qquad n\in\Nzero.
\end{equation}
These are the velocity, temporal acceleration, temporal jerk, and higher
Eulerian responses of the fluid.  They are read at a fixed spatial position.
The acceleration of a particle moving with the fluid is instead
\[
  D_tu=\partial_tu+(u\cdot\nabla)u.
\]
The transport term has not been discarded; the Eulerian Tower keeps it inside
the recurrence.  The Tower is not an extra model placed beside Navier--Stokes.
The original equation generates every rung.

The first row is simply the original equation solved for its time response:
\[
  V_1
  =\nu\Delta V_0-(V_0\cdot\nabla)V_0-\nabla Q_0.
\]
When time is differentiated once, the derivative may land on either factor in
the transport term:
\[
  V_2
  =\nu\Delta V_1
  -(V_1\cdot\nabla)V_0
  -(V_0\cdot\nabla)V_1
  -\nabla Q_1.
\]
At the next step, there are two ways to place one derivative on each transport
factor, and the middle term records both:
\[
  V_3
  =\nu\Delta V_2
  -(V_2\cdot\nabla)V_0
  -2(V_1\cdot\nabla)V_1
  -(V_0\cdot\nabla)V_2
  -\nabla Q_2.
\]
The coefficients are not decoration.  They count every way the \(n\) time
derivatives can be distributed across the velocity doing the transporting and
the velocity being transported.  The complete row is
\begin{equation}\label{eq:BodyTemporalVelocityRecurrence}
  V_{n+1}
  +\sum_{a=0}^{n}\binom{n}{a}(V_a\cdot\nabla)V_{n-a}
  +\nabla Q_n
  =\nu\Delta V_n,
  \qquad
  \nabla\cdot V_n=0.
\end{equation}

Pressure does not get to evolve independently beside these rows.
Incompressibility determines it from the velocity at the same time.  Taking
the divergence of the base equation and choosing the decaying whole-space
normalization gives
\begin{equation}\label{eq:PressureNormalization}
  -\Delta p=\partial_i\partial_j(u_i u_j),
  \qquad
  p=R_iR_j(u_i u_j),
\end{equation}
where \(R_j=\partial_j(-\Delta)^{-1/2}\).  After \(n\) time derivatives,
\begin{equation}\label{eq:BodyTemporalPressureRecurrence}
  Q_n
  =R_iR_j
  \sum_{a=0}^{n}\binom{n}{a}(V_a)_i(V_{n-a})_j.
\end{equation}
Thus the velocity rows determine \(Q_n\), and \(\nabla Q_n\) returns to
\eqref{eq:BodyTemporalVelocityRecurrence} as part of the very response those
rows generate.  Transport, pressure, incompressibility, and viscosity remain
one coupled event at every height of the Tower.

The temporal climb has now exposed a spatial pressure.  Every new response
contains \(\Delta V_n\); transport differentiates products of the earlier
responses; pressure is recovered from spatial derivatives of those products.
We entered the Tower because the hypothetical endpoint made the fluid respond
faster and faster.  The equation has answered by placing higher spatial order
inside every new response.  We return to the critical height to see exactly
what that spatial direction contains before naming the full two-index object.
